% !TEX root = ./ECEN-635 - Homework 3 - Brody Jordan.tex

\documentclass[11pt]{article}

\usepackage[
    assignmentDueDate={September 19, 2025 at 11:59 pm}
]{C:/Users/bljor/Sync/Courses/assignment}

\usepackage{bm}
\usepackage{siunitx}
% bold hat command
\newcommand{\bh}[1]{\mathbf{\hat{#1}}}

\begin{document}
\makeAssignmentTitle

% ————— PROBLEM 1 ————— %
\begin{homeworkProblem}
The field $\bm{\mathcal{E}}(t) = 5\sin(20 \times 10^9 t + kz) \mathbf{\hat{y}} ~ (\unit{V/m})$
describes a wave propagating in a medium which has the permittivity of free space with
$\varepsilon_r = 3$. Answer the following questions concerning this wave. Be sure to 
include dimensions in your answers:

\begin{enumerate}[label=(\alph*), topsep=5pt]
    \item What is the frequency $f$ of the wave?
    \item What is the wave amplitude?
    \item What is the propagation direction?
    \item What is the phase velocity of the wave $v_p$?
    \item What is the wavelength $\lambda$ of the wave?
    \item What is the intrinsic impedance of the wave $\eta$?
    \item What is the wave number (propagation constant) $k$?
\end{enumerate} \vspace{1\baselineskip}

\solution


\end{homeworkProblem}

\pagebreak

% ————— PROBLEM 2 ————— %
\begin{homeworkProblem}

In lossy insulating material a 2.45 GHz plane wave is propagating in a casserole in the $z$-
direction and polarized in the $y$-direction, with an amplitude of 20 V/m, a conductivity $\sigma = 2.0$ S/m, 
and a dielectric constant $\varepsilon_r = 4$. Find the lost tangent. Hint: this is a trivial question but the
loss tangent is of primary importance when ordering a material. \\

\solution


\end{homeworkProblem}

\pagebreak

% ————— PROBLEM 3 ————— %
\begin{homeworkProblem}

The first step in analytically solving for the fields in any situation in electromagnetics is to obtain the vector 
Helmholtz equation from Maxwell's equations. In class, starting with the Maxwell's equations the electric field 
vector Helmholtz equation was derived. Starting with Maxwell's equations derive the vector Helmholtz equation for 
the magnetic field. Write down and number Maxwell's equations and at each step in the derivation tell which equation 
you are going to use for the next step. If you use a vector identity, before using it write down the identity from 
those in the back cover of the book. Be sure to use the proper vector symbol and a line above the letter to indicate 
each vector. \\

\solution


\end{homeworkProblem}

\pagebreak

% ————— PROBLEM 4 ————— %
\begin{homeworkProblem}

In each case below (i) show mathematically that the field given is a straight line, a circle or an ellipse (ii) find 
an expression for the angle between the field vector and the x-axis. (iii) In each case describe the type of 
polarization. Hint: Remember the polarization is found by finding the phase of one of field components with respect 
to the other. A trick to find the difference in phase is to add or subtract a phase to both field components. By 
doing this you can cause one component to have zero phase and the other component will have the difference in the 
two phases. 

\begin{enumerate}[label=(\alph*), topsep=5pt]
    \item $\mathbf{E} = (j\bh{x} + \bh{y}) e^{-jkz}$
    \item $\mathbf{E} = [(1 + j)\bh{y} + (1-j) \bh{z}] e^{-jkx}$
    \item $\mathbf{E} = [(2 + j)\bh{x} + (3-j) \bh{z}] e^{-jky}$
    \item $\mathbf{E} = (j \bh{x} + j2\bh{y}) e^{-jkz}$
\end{enumerate} \vspace{1\baselineskip}

\solution


\end{homeworkProblem}

\pagebreak

% ————— PROBLEM 5 ————— %
\begin{homeworkProblem}
20s 322 hmwk 3 problem 3
\end{homeworkProblem}

\end{document}